1. What is a variable in JavaScript?
   >>> A variable in JavaScript is used to store a value.

2. What are the three keywords used to create variables?
    >>> Var, Let, const

3. Write the syntax to create a variable using var.
    >>> var age = 22;

4. Write the syntax to create a variable using let.
    >>> let name = "Unknown";

5. Write the syntax to create a variable using const.
    >>> const country = "India";

6. What is declaration?
    >>> Declaration means creating a variable by specifying its name, without assigning a value.
        var age
        let age
        const age

7. What is initialization?
    >>> Initialization means assigning a value to a variable when it is created.
        var age = 22;
        let age = 22;
        const age = 22;

8. What is reassignment?
    >>> Reassignment means assigning a new value to an already existing variable.
        var age = 22;
        age = 23;

9. What is redeclaration?
    >>> Redeclaration means declaring the same variable again in the same scope.
        var age = 22;
        var age = 23;

10. Which keyword allows redeclaration?
    >>>The var keyword allows redeclaration in the same scope.

11. Which keyword allows reassignment?
    >>> var and let keyword allows reassignment.

12. Which keyword requires initialization when declared?
    >>> const keyword requires initialization when declared.

13. Identify the declaration and initialization:
    let age = 25;
    >>> Declaration: let age
        Initialization: = 25

14. What is the value of a?
    var a = 100;
    console.log(a);
    >>> 100

15. Change the value of this variable to 200:
    let number = 100;
    >>> number = 200;
        console.log(number);
        Output = 200;

16. What will be the output?
    var a = 10;
    console.log(a);
    >>> Output = 10;

17. What will be the output?
    var a = 10;
    a = 20;
    console.log(a);
    >>> Output = 20;

18. What will be the output?
    var a = 10;
    var a = 30;
    console.log(a);
    >>> Output = 30;

19. Write a var variable named name with the value "John".
    >>> var name = "John";

20. Create a var variable named price with the value 500.
    >>> var price = 500;

21. Reassign price from 500 to 1000.
    >>> var price = 500;
        price = 1000;

22. What will be the output?
    var x = 50;
    x = 100;
    console.log(x);
    >>> Output = 100;

23. Can a var variable be reassigned?
    >>> Yes, var variable can be reassigned.

24. Can a var variable be redeclared?
    >>> Yes, var variable can be redeclared.

25. Write an example of var redeclaration.
    >>> var city = "Hyderabad";
        var city = "Bangalore";
        console.log(city);

26. Create a let variable named age with the value 25.
    >>> let age = 25;

27. What will be the output?
    let age = 20;
    age = 30;
    console.log(age);
    >>> output = 30;

28. Can a let variable be reassigned?
    >>> Yes, let variable can be reassigned.
29. Can a let variable be redeclared?
    >>> No, let variable cannot be redeclared in the same scope.
30. Find the error:
    let name = "John";
    let name = "David";
    >>>  Error: let cannot be redeclared in the same scope

31. Create a let variable called city and assign "Chennai".
    >>> let city = "Chennai";

32. Change the value of city to "Salem".
    >>> let city = "Chennai";
        city = "Salem";

33. What will be the output?
    let x = 10;
    x = 50;
    console.log(x);
    >>> 50;

34. Write a let variable called salary with the value 25000.
    >>> let salary = 25000;

35. Reassign salary to 30000.
    >>> let salary = 25000;
        salary = 30000;

36. Create a const variable called pi with the value 3.14.
    >>> const pi = 3.14;

37. Can a const variable be reassigned?
    >>> No, const variable cannot be reassigned.

38. Can a const variable be redeclared?
    >>> No, const variable cannot be redeclared in the same scope.

39. What is wrong with this code?
    const age;
    age = 25;
    >>>  const variable requires initialization at the time of declaration

40. What happens here?
    const price = 500;
    price = 1000;
    >>> Error - a const variable cannot be reassigned.

41. Create a const variable called country with the value "India".
    >>> const country = "India";

42. What will be the output?
    const x = 100;
    console.log(x);
    >>> 100;

43. Which keyword should you use if the value should not be reassigned?
    >>> const

44. What is the difference between let and const?
    >>> let - value can be changed.
        const - value cannot be changed.

45. What is the difference between var and const?
    >>> var - can be reassigned and redeclared.
        const - cannot be reassigned or redeclared.

46. Write JavaScript code to print Hello World using console.log().
    >>> const x = "Hello World";
        console.log(x);

47. Write JavaScript code to print the number 500 using console.log().
    >>> var x = 500;
        console.log(x);

48. What is the purpose of console.warn()?
    >>> console.warn() is used to display a warning message in the console.

49. What is the purpose of console.error()?
    >>> console.error() is used to display an error message in the console.

50. What is the purpose of each?
    alert() --->>> Shows a message in a popup box.
    prompt() --->>> Takes input from the user.
    confirm() --->>> Asks the user to confirm OK or Cancel.
    document.writeln() --->>> Writes content directly to the webpage.
    console.log() -->>> Prints information in the browser console.

<<<--- Practical Questions --->>>

i>Create a variable for student name, age, and mark and print all three.
    >>> let studentName = "Sai";
        let age = 25;
        let mark = 85;

        console.log(studentName);
        console.log(age);
        console.log(mark);

ii> Ask the user's name using prompt() and display it using alert().
    >>> let name = prompt("Enter your name:");
        alert(name);

iii> Ask the user's age using prompt() and print it using console.log().
    >>> let age = prompt("Enter your age:");
        console.log(age);

iv> Ask the user a question using confirm().
    >>> let answer = confirm("Are you learning JavaScript?");
        console.log(answer);

v> Ask the user's name and display it on the webpage using document.writeln().
    >>> let name = prompt("Enter your name:");
    document.writeln(name);

